\section*{Chapitre \ref{ch:s2:government-bond-relative-value} --- Valeur relative des obligations d'État}

\begin{solution}{exo:s2:government-bond-relative-value:1}
$\ln 0.5/\ln 0.5^{1/40} = 40$ jours.
\end{solution}

\begin{solution}{exo:s2:government-bond-relative-value:2}
$2 \times 4.99 - 4.85 - 5.11 = 0.02$ points de pourcentage~: 2.0 points de base.
\end{solution}

\begin{solution}{exo:s2:government-bond-relative-value:3}
Les moindres carrés rendent les résidus orthogonaux aux régresseurs~: $X^\top(y - X\hat\beta) = 0$. Un portefeuille dont les poids sont proportionnels aux résidus
est déjà neutre aux facteurs ajustés~; la couverture ne compte qu'une fois que le portefeuille ne sélectionne qu'une partie des obligations.
\end{solution}

\begin{solution}{exo:s2:government-bond-relative-value:4}
Le bruit ajoute de la variance au résidu d'aujourd'hui sans aucune persistance, si bien que la régression de demain sur aujourd'hui est biaisée vers zéro~: la
persistance ajustée est trop faible et la demi-vie trop courte.
\end{solution}

\begin{solution}{exo:s2:government-bond-relative-value:5}
Elles se négocient souvent avec une prime pour leur liquidité et leur specialness en repo, ce qui fausserait une courbe construite sur des obligations ordinaires. La prime
suggère de vendre l'émission on-the-run contre une émission off-the-run voisine, si son coût en repo peut être compensé par le carry.
\end{solution}

\begin{solution}{exo:s2:government-bond-relative-value:6}
Retirer trois facteurs élimine presque tout le risque de courbe, si bien que la volatilité brute baisse le plus~; mais la couverture ajoute de petites positions sur les quarante obligations qui
changent chaque semaine, et à 0.25 point de base, le surcroît de rotation coûte plus que ne vaut le risque qu'il retire.
\end{solution}

\begin{solution}{exo:s2:government-bond-relative-value:7}
Le portefeuille neutre en DV01 passe de 1.13 à 2.54 et le portefeuille neutre aux facteurs de 0.69 à 3.48~: avec une exécution bon marché, la meilleure couverture l'emporte. Cette stratégie revient
à celui qui négocie au meilleur coût, souvent un dealer ou un fonds dont les coûts d'exécution et de financement sont faibles.
\end{solution}

\begin{solution}{exo:s2:government-bond-relative-value:8}
Les 18 jours viennent d'une régression à un jour biaisée par le bruit de cotation et par la partie rapide de l'erreur~; à 120 jours, 40~\% d'un résidu subsiste encore.
Clôturer après un mois laisse non capturée l'essentiel du retour à la moyenne attendu et oblige à payer de nouveau la rotation.
\end{solution}

\begin{solution}{pb:s2:government-bond-relative-value:1}
\begin{enumerate}
\item La courbe quotidienne du Trésor américain publiée par le Board depuis 1961, ajustée avec l'extension de Svensson au modèle de Nelson--Siegel.
\item Obligations remboursables par anticipation et flower bonds, titres à moins de trois mois de l'échéance, bons du Trésor, obligations à vingt ans depuis 1996, et les premières
émissions off-the-run~: elles portent des clauses, une segmentation ou des primes de liquidité que les obligations ordinaires n'ont pas.
\item 87.9~\%, 8.2~\% et 2.0~\%.
\item Rendement de marché moins le rendement de la courbe ajustée à l'échéance de l'obligation.
\item Milieu contre ailes, neutre en DV01 et en pente~; moyenne de $-1.6$ points de base, écart-type 26.6, demi-vie 105 jours.
\item Un classement des obligations selon leur résidu pour choisir lesquelles acheter et vendre.
\item 2 points de base qui reviennent en 40 jours, 2 en 250 jours, vus à travers 0.5 point de base de bruit.
\item Cinq obligations les moins chères en position longue, cinq les plus chères en position courte, chaque semaine~; non couvert (longue seulement), neutre en DV01 ou neutre sur trois facteurs.
\item 0.36, 3.46 et 5.41 brut~; 0.25, 1.13 et 0.69 net.
\item C'est surtout une position de duration.
\item À 0.25 point de base, la rotation de la couverture par facteurs coûte plus que ne vaut la réduction de risque qu'elle apporte~; à 0.1 point de base, ce n'est pas le cas.
\item 81~\% restent après 20 jours, 59~\% après 60, 40~\% après 120.
\item \textbf{Résultat nommé.} Net des coûts, le portefeuille neutre en DV01 affiche un ratio de Sharpe de 1.13 et le portefeuille neutre aux facteurs 0.69 (5.41 brut)~; le
demi-vie à un jour des résidus est de 17.9 jours, mais 40~\% d'un résidu subsiste après 120 jours.
\item Le bruit de cotation et la partie lente de l'erreur lui échappent tous deux.
\item Financement~: décotes, roll de repo, specialness, positions encombrées.
\item Fourchette, repo et rotation de la couverture.
\item Ajuster sur les cotations connues à l'époque, négocier au cours acheteur et au cours vendeur, inclure le repo, mesurer la décroissance à l'horizon de détention.
\item Off-the-run contre on-the-run.
\item Le chapitre~11 négocie une obligation contre son contrat à terme, financée en repo à fort effet de levier.
\item Liquidité et bilan pour les investisseurs dont les demandes éloignent les obligations de la courbe.
\end{enumerate}
\end{solution}

\begin{solution}{iq:s2:government-bond-relative-value:1}
Son rendement comparé à une courbe ajustée sur des obligations comparables~: l'offre, la demande liée aux indices et la demande réglementaire, la liquidité, la specialness en repo, les effets de coupon et de fiscalité font
que les obligations se situent au-dessus ou au-dessous.
\end{solution}

\begin{solution}{iq:s2:government-bond-relative-value:2}
Une forme paramétrique lisse (Nelson--Siegel--Svensson) ou une spline ajustée sur des titres anciens et des obligations pondérées par la duration~; on exclut les bons du Trésor, les obligations proches
de l'échéance, les émissions on-the-run et les obligations à caractéristiques spécifiques.
\end{solution}

\begin{solution}{iq:s2:government-bond-relative-value:3}
Résoudre les poids des ailes pour que le DV01 total et l'exposition à un mouvement de pente soient nuls~; la position parie toujours sur la courbure au milieu, sur les trois
obligations, sur leurs propres résidus et sur leurs taux de repo.
\end{solution}

\begin{solution}{iq:s2:government-bond-relative-value:4}
Les décotes montent et forcent des ventes, le repo ne peut plus être roulé, les résidus s'élargissent avant de revenir, de nombreux fonds détiennent les mêmes positions et vendent ensemble.
\end{solution}

\begin{solution}{iq:s2:government-bond-relative-value:5}
Entrées~: prix pied de coupon ou rendements des obligations éligibles~; exclusions selon une règle~; un ajustement aux moindres carrés pondérés~; sorties~: paramètres, rendements ajustés, résidus
et leur historique~; contrôles~: erreur d'ajustement, sauts de paramètres, cours périmés, nouvelles émissions.
\end{solution}

\begin{solution}{iq:s2:government-bond-relative-value:6}
La pente est $\mathrm{Cov}(e_{t+1} + \eta_{t+1}, e_t + \eta_t)/\mathrm{Var}(e_t + \eta_t) = \phi\sigma_e^2/(\sigma_e^2 + \sigma_\eta^2)$. Elle est inférieure
à $\phi$, donc $\ln 0.5/\ln(\text{pente})$ sous-estime la demi-vie~; le biais croît avec la part du bruit dans la variance du résidu.
\end{solution}
